\chapter*{Preface}
\addcontentsline{toc}{chapter}{Preface}

Book 1's Part II introduced, in the space of two chapters, the idea that entropy behaves as an ideal within a ring built to support the Lie-algebraic commutator $[\PhiField, \SEntropy]$. That introduction was necessarily brief, since Book 1's ring $\mathcal{R} = \langle \PhiField, \vField, \SEntropy \rangle$ is built around transport by the flow field $\vField$, and the ideal-theoretic content of entropy was only ever a supporting player in an argument whose main business was curvature. The present book removes that supporting role and makes entropy-as-ideal the whole subject, on a ring built without $\vField$ at all: a purely commutative ring of scalar and entropy observables, in which the interesting structure is not how potential and entropy fail to commute under transport, but how the ideal generated by entropy organizes the ring's own spectrum of possible states.

This is not a difference of emphasis but a difference of mathematical object. Book 1's ring is noncommutative because transport, in general, does not commute with itself along different paths; the present ring is commutative because, once flow is set aside, multiplying scalar and entropy observables in either order describes the same static configuration. The two rings are related, as Chapter 11 below shows, but neither is a restriction or an extension of the other in any simple sense, and the present book develops its own ring on its own terms before returning, in its closing part, to say precisely how the two fit together.

\chapter*{Concept Map}
\addcontentsline{toc}{chapter}{Concept Map}

Three constructions organize this book. The ring $\mathcal{R}_2$, generated by $\PhiField$ and $\SEntropy$ alone under ordinary commutative multiplication, replaces Book 1's three-generator, flow-mediated ring as the object of study. The maximal ideal $\mathfrak{m}_{\SEntropy}$, previewed briefly in Book 1, is here given the actual algebraic condition for its maximality and used to build a genuine quotient ring and spectrum, rather than asserted as a bare fact. The spectrum $\mathrm{Spec}\,\mathcal{R}_2$, the space of prime ideals of this ring, is developed as an object in its own right, with physical and ethical order read as properties of specific localizations within it. The closing part returns to Book 1 and asks how this static, commutative picture relates to that book's dynamical, noncommutative one.

\chapter*{Reading Note}
\addcontentsline{toc}{chapter}{Reading Note}

Where Book 1 asked the reader to think dynamically, in terms of transport and commutators, this book asks for a different habit of mind: think statically, in terms of what can be said about a ring's structure without reference to how its elements evolve. The two habits are not in tension, and Part VI shows how to move between them, but conflating them prematurely is exactly the mistake an earlier working draft of this corpus made, and the present book is deliberately paced to avoid repeating it.

\part{The Ring of Scalar--Entropy Observables}

